\documentclass[12pt]{article}
\usepackage[T1]{fontenc}
\usepackage[utf8]{inputenc}
\usepackage{lmodern}
\usepackage{textcomp}
\usepackage{enumitem}
\usepackage{geometry}
\geometry{margin=1in}
\begin{document}
\begin{center}
Answer Key - Constitutional Law - Undergraduate Level Assessment 5 \
\small (Answers are ordered exactly as in the question paper.)
\end{center}

\section*{Multiple Choice}
\begin{enumerate}[label=\arabic*.]
\item \textbf{Answer:} B
\\ \textit{Options:} A) They are acts that States perform as practice in the context of custom; B) They are acts creating unilateral legal obligations to the acting State; C) Unilateral acts are simply political acts of State devoid of any legal effect; D) Unilateral acts are those that State perform in order to be bound by a treaty
\\ \textit{Note:} Unilateral acts refer to actions taken by a State that create binding legal obligations without the need for agreement from other parties, thereby establishing certain legal responsibilities solely based on the will of the acting State.
\item \textbf{Answer:} C
\\ \textit{Options:} A) Because Aristotle justified slavery.; B) Because Kant failed to distinguish individual from social morality?; C) Because Aristotle believed that man is a 'social animal'; D) Because Kant regarded the individual as unimportant.
\\ \textit{Note:} The correct answer highlights that communitarianism aligns with Aristotle's philosophy because it emphasizes the importance of community and social relationships, reflecting Aristotle's view of humans as inherently social beings who find their identity and purpose within the context of society.
\item \textbf{Answer:} D
\\ \textit{Options:} A) You must not consciously harm another person.; B) You must always act in the best interests of the community.; C) You must treat human beings as means rather than ends.; D) You must act according as if your values apply to everyone.
\\ \textit{Note:} Kant's 'categorical imperative' posits that one should act according to maxims that can be universally applied, meaning that your actions should reflect principles that could be accepted as universal laws, thus aligning personal values with a broader ethical framework.
\item \textbf{Answer:} A
\\ \textit{Options:} A) It is misinterpreted as a prediction.; B) His concept of status is misrepresented.; C) It is taken literally.; D) His idea is considered inapplicable to Western legal systems.
\\ \textit{Note:} The correct answer is based on the fact that the aphorism is often misconstrued as a forward-looking prophecy about societal evolution rather than a descriptive observation of historical trends in legal and social relationships.
\item \textbf{Answer:} A
\\ \textit{Options:} A) It regards morals and law as inseparable.; B) It perceives law as commands.; C) It regards a legal order as a closed logical system.; D) It espouses the view that there is no necessary connection between morality and law.
\\ \textit{Note:} The correct answer is accurate because legal positivism fundamentally asserts that law and morality are distinct and that the validity of a law is not dependent on its moral content, contrary to the claim that they are inseparable.
\end{enumerate}

\section*{True / False}
\begin{enumerate}[label=\arabic*.]
\setcounter{enumi}{5}
\item \textbf{Answer:} True
\\ \textit{Options:} A) True; B) False
\\ \textit{Note:} Austin's theory of jurisprudence emphasizes that positive law, which is law created and enforced by a recognized authority, is the primary subject of legal study, distinguishing it from moral or natural law.
\item \textbf{Answer:} True
\\ \textit{Options:} A) True; B) False
\\ \textit{Note:} The Universal Declaration of Human Rights (UDHR) is indeed a resolution adopted by the UN General Assembly on December 10, 1948, establishing fundamental human rights that are universally protected.
\item \textbf{Answer:} False
\\ \textit{Options:} A) True; B) False
\\ \textit{Note:} The protective principle of jurisdiction allows a state to exercise authority over its nationals for offenses committed abroad, but it also considers the connection between the offense and the state, including the harm caused to its interests.
\item \textbf{Answer:} True
\\ \textit{Options:} A) True; B) False
\\ \textit{Note:} The statement is true because the doctrine of stare decisis establishes that courts are bound to follow the legal principles established in previous cases to ensure consistency and predictability in the law.
\item \textbf{Answer:} True
\\ \textit{Options:} A) True; B) False
\\ \textit{Note:} Individuals possess international legal personality, allowing them to have rights and obligations under international law, but their capacity to act and engage in legal relations is significantly more restricted than that of states and international organizations, which have broader recognition and authority in international affairs.
\end{enumerate}

\section*{Long Form Response}
\begin{enumerate}[label=\arabic*.]
\setcounter{enumi}{10}
\item \textbf{Answer:} The Social Fact Thesis, as articulated by Ken Himma within the framework of positivism, posits that law is fundamentally a social construction or artifact, shaped by the collective actions and agreements of society. This perspective emphasizes that legal norms and rules emerge not from inherent moral principles but from social practices and conventions agreed upon by a community. As a result, constitutional law, which governs the structure and operation of government, can be understood as a product of societal consensus reflecting the values and interests of the people. This view implies that constitutional law is subject to change and evolution as societal norms shift, highlighting the dynamic relationship between law and society. Consequently, it reinforces the idea that legal systems must be responsive to the needs and beliefs of the communities they serve.
\\ \textit{Note:} correct because it accurately captures the essence of the Social Fact Thesis by Ken Himma, which asserts that law is derived from social agreements rather than moral truths, thereby framing constitutional law as a product of societal consensus and practices.
\item \textbf{Answer:} The 'fair play' argument posits that individuals have a moral duty to obey the law because they benefit from the social order and protections that the legal system provides. However, this argument is challenged by the claim that the legal system is inherently unfair, as it often perpetuates inequality and discrimination against certain groups. For instance, systemic biases can lead to disproportionate penalties for marginalized communities, undermining the notion of mutual benefit that the fair play argument relies on. Consequently, if the legal system is perceived as biased or unjust, individuals may justifiably question their obligation to obey laws that do not apply equitably to all. This disconnect between the ideal of fair play and the reality of legal inequities highlights a significant challenge to the moral foundation of the duty to obey the law.
\\ \textit{Note:} The answer correctly identifies the 'fair play' argument as a justification for obeying the law based on the benefits received from the legal system, while simultaneously addressing the counterargument that inherent unfairness within the system undermines the moral obligation to comply.
\end{enumerate}

\vfill
\noindent\textit{End of Answer Key}
\end{document}